\documentclass[10pt,a4paper]{amsart}
\usepackage[margin=19mm]{geometry}
\usepackage{amsmath,amssymb,mathtools}
\usepackage[scr=rsfs,cal=euler]{mathalfa}
\usepackage{xfrac}
\usepackage[unicode,hidelinks,bookmarks=false]{hyperref}
\newcommand{\ie}{i.e.\ }
\newcommand{\cf}{cf.\ }
\newcommand{\resp}{respectively\ }
\newcommand{\Subsection}{\subsection}
\providecommand{\nobreakdash}{\nobreak\hbox{-}}
\setlength{\emergencystretch}{2em}
\multlinegap0pt
\usepackage{bm}
\usepackage{bbm}
\usepackage{dsfont}
\usepackage{xspace}
\usepackage{cancel}
\usepackage{mathtools}
\usepackage{cleveref}
\usepackage{amsthm}
\makeatletter
\@ifundefined{theorem}{\newtheorem{theorem}{Theorem}}{}
\@ifundefined{lemma}{\newtheorem{lemma}[theorem]{Lemma}}{}
\@ifundefined{prop}{\newtheorem{prop}[theorem]{Proposition}}{}
\makeatother
\DeclareMathOperator{\Ent}{Ent}
\DeclareMathOperator{\TBM}{\mathsf{TBM}}
\DeclareMathOperator{\TMCP}{\mathsf{TMCP}}
\DeclareMathOperator{\diam}{diam}
\DeclareMathOperator{\spt}{spt}
\newcommand{\N}{\mathbb{N}}
\newcommand{\Pcal}{\mathcal{P}}
\newcommand{\R}{\mathbb{R}}
\newcommand{\m}{\mathfrak{m}}
\newcommand{\narrowto}{\rightharpoonup}
\renewcommand{\epsilon}{\varepsilon}
\title[The equivalence between timelike Ricci curvature and the tim]{The equivalence between timelike Ricci curvature and the timelike Brunn Minkowski inequality on synthetic Lorentzian spaces}
\author{Osama Farooqui}
\date{Source: arXiv:2604.11766v1 (CC BY 4.0)}
\begin{document}
\maketitle

\noindent\emph{Source and licence.} The equivalence between timelike Ricci curvature and the timelike Brunn Minkowski inequality on synthetic Lorentzian spaces. Version arXiv:2604.11766v1 (CC BY 4.0), distributed under the Creative Commons Attribution 4.0 International licence, as recorded in the arXiv licence metadata for this exact version and shown on the abstract page https://arxiv.org/abs/2604.11766v1. Full rights notice: licenses/arxiv-2604.11766-CC-BY-4.0.txt.\par\smallskip

\noindent\textit{Editorial scope.} Counted unit: the Theorem whose source label is \texttt{TBMtoTMCP} in the article (source0.tex lines 808--809, source label \texttt{TBMtoTMCP}), delivered together with the article's own complete original proof of it (source0.tex lines 810--871, one \texttt{proof} environment of 62 lines). No mathematics is written, repaired or re-derived: statement and proof are the article's own text line for line. Required context retained uncounted: simplesequence (lines 635--643); midpointapproxlemma (lines 683--688); midpointTMCP (lines 747--753); DiracTBM (lines 774--776). Every definition, display and label of those lines is retained unchanged. Omitted from this package and not redistributed: 1--634, 644--682, 689--746, 754--773, 777--807, 872--1111. Nothing inside a retained excerpt is omitted or altered: every retained body is delivered exactly as the article's lines are written, so no omission receipt of any kind accompanies this package. Citations: the retained excerpts carry no citation command at all, so no bibliography entry of the article is transcribed and no reference is redistributed. Reference adaptation: none was needed and none was made; every reference inside the delivered unit resolves inside it, so no unresolved reference or citation is produced. Printed slips kept literal: no printed word, symbol, hypothesis or number of the counted statement or of its proof was altered. Layout and class: the article's own class is \texttt{article} with packages including fontenc, inputenc, babel, amsmath, amsfonts, amssymb, amsthm, diffcoeff, geometry, thmtools, hyperref, cleveref, biblatex; the wrapper is the lane's pinned \texttt{amsart} editorial wrapper at 10pt on A4, which restates the theorem-like environments the retained text needs and adds the article's own macro definitions that the retained text needs (\texttt{\textbackslash Ent}, \texttt{\textbackslash N}, \texttt{\textbackslash Pcal}, \texttt{\textbackslash R}, \texttt{\textbackslash TBM}, \texttt{\textbackslash TMCP}, \texttt{\textbackslash diam}, \texttt{\textbackslash epsilon}, \texttt{\textbackslash m}, \texttt{\textbackslash narrowto}, \texttt{\textbackslash spt}), with extra wrapper packages (bm, bbm, dsfont, xspace, cancel, mathtools, cleveref). The delivered numbering is the unit's own. Assets: no retained span contains a figure environment, a graphics-inclusion command, a PDF-page inclusion, a table or a TikZ picture; no image file is copied into this package, the package declares no asset dependency and no image file is redistributed with it. Licence: version arXiv:2604.11766v1 of this article is distributed under the Creative Commons Attribution 4.0 International licence, as recorded in the arXiv licence metadata for this exact version and shown on the abstract page https://arxiv.org/abs/2604.11766v1. A full rights notice is delivered at licenses/arxiv-2604.11766-CC-BY-4.0.txt. Characters: every retained excerpt is pure ASCII, so the delivered engine, XeLaTeX with its default Latin Modern text font, reports no missing character.
\begin{prop}[Existence of simple sequences]\label{simplesequence} Let $\mu\in \Pcal_c^{ac}(X)$. Then there exists a sequence $\{\mu^n\}_{n\in\N}$ of simple measures, which we call a ``simple sequence", which satisfies the following properties: 
\begin{enumerate}
\item (Narrow Convergence). As $n\to\infty$, $\mu^n\narrowto\mu$. 
\item (Controlled Decomposition). For each metrizer $d$ of $X$, and each positive sequence $\delta_n\to 0$, we may write $\mu^n:= \sum_{j=1}^{M_n}\lambda_j^n \m_{A_j^n}$, where for each $n\in\N$ and each $j=1,\dots, M_n$, $A_j^n$ is compact and contained in $\spt\mu_0$, and $\diam_d(A_j^n)<\delta_n$. 

\item (Continuity along Entropy) For each $N\in\N$, $U_N(\mu^n)\to U_N(\mu)$, and $S_N(\mu^n)\to S_N(\mu)$. 

\end{enumerate}
\end{prop}

\begin{lemma}[Midpoint approximation lemma]\label{midpointapproxlemma} Let $(X,\ell, \tau,\m)$ be a compact measured synthetic spacetime, and suppose $\ell_+$ is continuous. Then for each metrizer $d$ of $X$ and each $\epsilon>0$, there exists $\delta>0$ such that for each pair of totally timelike measures $(\mu_0,\mu_1)\in\Pcal(X)^2$ and each $\mu_t\in \Pcal(X)$ such that $\spt\mu_t\subset G_t(\spt\mu_0,\spt\mu_1)$, if $\diam_d(\spt\mu_0),\diam_d(\spt\mu_1)<\delta$ then we have 
\begin{align}
\ell_q^q(\mu_0, \mu_t)\geq t^q\ell^q_q(\mu_0,\mu_1)-\epsilon;\quad 
\ell_q^q(\mu_t, \mu_1)\geq (1-t)^q\ell^q_q(\mu_0,\mu_1)-\epsilon.
\end{align}
\end{lemma}

\begin{lemma}[$\TMCP$ on $t$-midpoints]\label{midpointTMCP}
Let $X$ be a measured globally hyperbolic synthetic spacetime. Suppose there exists $K\in\R$ and $N>1$ such that, for every $\mu=\rho\m\in \Pcal_{c}^{ac}(X)$, every $x_0\in X$ such that $\spt\mu_0\subset I^-(x_0)$, and every $t\in (0,1)$, there exists a $t$-midpoint $\mu_{t}$ between $\mu$ and $\delta_{x_0}$ such that 
\begin{align}\label{midpointtmcp}
U_N(\mu_{t})\geq \sigma_{K/N}^{(1-t)}(\Lambda)U_N(\mu),
\end{align}
where $\Lambda:=\sqrt{\int \ell^2(x,x_0)d\mu(x)}$. Then $X\in \TMCP^e(K,N)$. 
\end{lemma}

\begin{lemma}[$\TBM(K,N)$ for Dirac targets]\label{DiracTBM} Let $(X,\ell,\tau,\m)$ be a measured synthetic spacetime, with $\ell_+$ continuous. Suppose $X\in \TBM^*(K,N)$. Let $x_0\in X$ and let $A\subset X$ be a compact set of positive measure such that $A\subset I^-(x_0)$. Then
\[\m^{1/N}(G_t(A,x_0))\geq \sigma_{K/N}^{(1-t)}(\Theta(A,x_0))\m^{1/N}(A).\]
\end{lemma}

\begin{theorem}[$\TBM^*(K,N)$ spaces are $\TMCP^e(K,N)$ spaces]\label{TBMtoTMCP} Fix $0<q<1$. Let $(X,\ell,\tau,\m)$ be a timelike $q$-essentially non-branching measured globally hyperbolic synthetic spacetime, with $\ell_+$ continuous. If $X\in \TBM^*(K,N)$, then $X\in \TMCP^e(K,N)$. 
\end{theorem}

\begin{proof}
Let $\mu\in\Pcal_c^{ac}(X)$, and $x_0\in X$ be such that $(\spt\mu,x_0)$ is timelike. In light of \Cref{midpointTMCP}, we aim to show that for each $t\in(0,1)$, there is a $t$-midpoint $\mu_t$ joining $\mu$ to $\delta_{x_0}$ such that 
\begin{align}\label{convexityinequality1}
U_N(\mu_t)\geq\sigma_{K/N}^{(1-t)}(\Lambda)U_N(\mu),
\end{align}
where $\Lambda:= \sqrt{\int \ell^2(x,x_0)d\mu(x)}$. Fix $t\in (0,1)$.


\paragraph{Step 1: Obtaining the $t$-midpoint $\mu_t$.}

From global hyperbolicity, the set $Y:=J(\spt\mu, x_0)$ is a compact synthetic spacetime. For each $n\in\N$, we apply \Cref{midpointapproxlemma} with the choice of metrizer $d$ of $Y$ and $\epsilon=1/n$ to obtain some $\delta_n>0$ such that for any pair $(\mu_0, \mu_1)\in \Pcal(Y)^2$ of compactly supported totally timelike measures, and any measure $\mu_t\in \Pcal(Y)$ with $\spt\mu_t\subset G_t(\mu_0, \mu_1)$,  if $\diam_d(\spt\mu_i)<\delta_n$ for $i=0,1$, then 
\[\ell_q(\mu_0, \mu_t) \geq t\ell_q(\mu_0, \mu_1)-\frac 1n; \quad \ell_q(\mu_t, \mu_1) \geq (1-t)\ell_q(\mu_0, \mu_1)-\frac1n.\]

WLOG we may assume $\delta_n\to 0$ as $n\to\infty$. By \Cref{simplesequence} there exists a simple sequence $\{\mu^n\}_{n\in\N}$ which converges to $\mu$. From the controlled decomposition property, we may write $\mu^n:=\sum_{j=1}^{M_n}\lambda_j^n \m_{A_j^n}\in \Pcal(Y)$, where for each $n$, the collection $\{A_j^n\}_{j=1}^{M_n}$ are compact, mutually disjoint, have finite positive $\m$-measure, $A_j^n\subset \spt\mu\subset Y$, and $\diam_d(A_j^n)<\delta_n$. 

Let $A_{t,j}^n:=G_t(A_j^n,x_0)\subset Y$. From \Cref{DiracTBM}, we know that $A_{t,j}^n$ has positive $\m$-measure. Since $A_{t,j}^n$ is compact and $\m$ is Radon, we also know that it has finite $\m$-measure. We define 
\[\mu_{t}^n:=\sum_{j=1}^{M_n}\lambda_j^n \m_{A_{t,j}^n}.\]

Observe that $\mu_t^n\in \Pcal(Y)$, and so $\mu_t^n$ admits a narrow subsequential limit $\mu_t$. Concavity of $\ell_q^q$ and \Cref{midpointapproxlemma} applied to the triple $(\mu_0,\mu_t,\mu_1):= (\m_{A_j^n}, \m_{A_{t,j}^n}, \delta_{x_0})$ then gives that
\begin{align*}
\ell_q^q(\mu^n,\mu_t^n) \geq \sum_{j=1}^{M_n}\lambda_j^n \ell_q^q(\m_{A_j^n}, \m_{A_{t,j}^n})
&\geq \sum_{j=1}^{M_n} \lambda_j^n t^q\ell_q^q(\m_{A_j^n}, \delta_{x_0})-\frac1n = t^q\sum_{j=1}^{M_n}\lambda_j^n \int\ell^q(x,x_0)d\m_{A_j^n}(x)-\frac 1n \\
& = t^q\int\ell^q(x,x_0) d\mu^n(x)-\frac 1n 
\end{align*}
Narrow convergence of $\mu^n\narrowto\mu$, continuity of $\ell_+^q$, and positivity of $\ell^q$ on $\spt\mu\times \{x_0\}$, and upper semicontinuity of $\ell_q$ then gives that, after passing to subsequence, that for all $\epsilon>0$,
\begin{align*}
\ell_q^q(\mu, \mu_t) \geq \limsup_{n\to\infty} \ell_q^q(\mu^n, \mu_t^n) = t^q\limsup_{n\to\infty}\int \ell^q(x,x_0) d\mu^n(x)-\frac 1n=t^q\int \ell^q(x,x_0)d\mu(x).
\end{align*}
Hence we may conclude that 
\[\ell_q(\mu,\mu_t)\geq t\ell_q(\mu,\delta_{x_0}).\]
A symmetric argument gives also that 
\[\ell_q(\mu_t,\delta_{x_0})\geq (1-t)\ell_q(\mu,\delta_{x_0}),\]
which proves that $\mu_t$ is indeed a $t$-midpoint between $\mu$ and $\delta_{x_0}$. \\

\paragraph{Step 2: Entropic Convexity along $\mu_t$.}


We now verify the entropic convexity inequality \eqref{convexityinequality1} along $\mu_t$. The $q$-essential non-branching assumption guarantees that for each $n$, the sets $A_{t,j}^n$ are mutually disjoint for all $j=1,\dots M_n$, since the $A_j^n$ are mutually disjoint. As a result, we have 

\begin{align*}
\Ent(\mu_t^n) &= \sum_{j=1}^{M_n}\Ent(\lambda_j^n\m_{A_{t,j}^n}) = \sum_{j=1}^{M_n} \lambda_j^n\log\frac{\lambda_j^n}{\m(A_{t,j}^n)}.
\end{align*}
Therefore we have 
\begin{align*}
U_N(\mu_t^n) &=\exp\left(-\frac1N\Ent(\mu_t^n)\right) =\prod_{j=1}^{M_n} \left(\frac{\m(A_{t,j}^n)}{\lambda_j^n}\right)^{\lambda_j^n/N} \geq \prod_{j=1}^{M_n}\left(\sigma_{K/N}^{(1-t)}(\Theta_j^n)\frac{\m^{1/N}(A_{0,j}^n)}{(\lambda_j^n)^{1/N}}\right)^{\lambda_j^n } \\
&= \prod_{j=1}^{M_n}\left(\sigma_{K/N}^{(1-t)}(\Theta_j^n)\right)^{\lambda_j^n} U_N(\mu^n),
\end{align*}
where $\Theta_j^n=\Theta(A_j^n,x_0)$. Now observe that log-convexity of $\theta\mapsto \sigma_{k}^{(t)}(\sqrt \theta)$ implies that 
\begin{align}\label{distortionestimate1}
\prod_{j=1}^{M_n} \sigma_{K/N}^{(1-t)}(\Theta_j^n)^{\lambda_j^n}\geq \sigma_{K/N}^{(1-t)}\left(\Big(\sum_{j=1}^{M_n} \lambda_j^n (\Theta_j^n)^2\Big)^{1/2}\right).
\end{align}
Continuity of $\ell$, compactness of $A_j^n$, and the fact that $\diam_d(A_j^n)<\delta_n$ gives that for each there exists $n_0\in\N$ such that for all $n>n_0$, and all $x_j^n\in A_j^n$, 
\[|\Theta_j^n-\ell(x_j^n,x_0)|<1/n.\]
It follows that, with $\Lambda^n:= \sqrt{\int \ell^2(x,x_0)d\mu^n}$, we have 
\[\left| \Big(\sum_{j=1}^{M_n} \lambda_j^n (\Theta_j^n)^2\Big)^{1/2}-\Lambda^n\right|<\frac 1n,\]
And so by monotonicity of $\theta\mapsto \sigma_{K/N}^{(t)}(\theta)$ we obtain 
\[U_N(\mu_t^n) \geq \sigma_{K/N}^{(1-t)}\left(\Lambda^n\mp \frac 1n\right)U_N(\mu^n), \]
where we choose $-\frac 1n$ if $K\geq 0$ and $+\frac 1n$ if $K<0$. Narrow convergence of $\mu^n$ to $\mu$ guarantees that $\Lambda^n\to \Lambda$, and the continuity along entropy guarantees that $U_N(\mu^n)\to U_n(\mu)$. Thus from upper semicontinuity we obtain
\[U_N(\mu_t)\geq \limsup_{n\to\infty} U_N(\mu_t^n) \geq \limsup_{n\to\infty}\sigma_{K/N}^{(1-t)}\left(\Lambda^n\mp \frac1n\right)U_N(\mu^n)= \sigma_{K/N}^{(1-t)}(\Lambda)U_N(\mu).\]
The last equality holds since $\sigma_{K/N}^{(1-t)}(\Lambda)>0$. 

\end{proof}

\end{document}
